%=============================================================================
% Section 6: Execution Layer Position in AI Native Stack
%=============================================================================
\section{Execution Layer Position in AI Native Stack}

\subsection{Five-Layer Architecture}

The execution layer sits within a broader AI Native implementation stack:

%-----------------------------------------------------------------------------
% Figure 2: Five-Layer Architecture
%-----------------------------------------------------------------------------
\begin{figure}[htbp]
\centering
\begin{tikzpicture}[
    layerbox/.style={
        draw,
        rectangle,
        minimum width=8.5cm,
        minimum height=0.75cm,
        align=center,
        font=\small
    },
    subbox/.style={
        draw,
        rectangle,
        minimum width=7cm,
        minimum height=0.65cm,
        align=center,
        font=\footnotesize
    },
    annot/.style={
        font=\footnotesize,
        text=gray!70!black
    },
    arrow/.style={
        -{Stealth[length=2.5mm, width=1.8mm]},
        thick,
        gray!60!black
    }
]

% User Interface
\node[layerbox, fill=gray!12] (ui) at (0, 7) {User Interface};

\draw[arrow] (0, 6.5) -- (0, 6.1);

% AI Native Applications
\node[layerbox, fill=blue!8] (apps) at (0, 5.5) {AI Native Applications};
\node[annot] at (0, 4.9) {(Agents, Workflows, Apps)};

\draw[arrow] (0, 4.5) -- (0, 4.1);

% AI Cognition Layer
\node[layerbox, fill=green!8] (cog) at (0, 3.5) {AI Cognition Layer};
\node[annot] at (0, 2.9) {(LLMs, Reasoning, Planning)};

\draw[arrow] (0, 2.5) -- (0, 2.1);

% AI Action Layer (container)
\node[draw, rectangle, minimum width=8.5cm, minimum height=2.8cm, fill=yellow!6] (action) at (0, 0.3) {};
\node[font=\small\bfseries, anchor=north west] at (-3.8, 1.6) {AI Action Layer};

% Orchestration/Policy subbox
\node[subbox, fill=green!12] (orch) at (0, 0.9) {Orchestration/Policy};

\draw[arrow] (0, 0.45) -- (0, 0.05);

% Execution Runtime subbox
\node[subbox, fill=orange!18] (exec) at (0, -0.6) {Execution Runtime};

\draw[arrow] (0, -1.1) -- (0, -1.5);

% Capability Ecosystem
\node[layerbox, fill=purple!8] (cap) at (0, -2.1) {Capability Ecosystem};
\node[annot] at (0, -2.75) {(Models, APIs, Tools, Devices)};

\draw[arrow] (0, -3.1) -- (0, -3.5);

% Infrastructure
\node[layerbox, fill=gray!15] (infra) at (0, -4.1) {Infrastructure};
\node[annot] at (0, -4.75) {(Cloud, Edge, Networks)};

\end{tikzpicture}
\caption{Five-layer AI Native implementation stack showing the position of the AI Execution Layer within the broader architecture.}
\label{fig:five-layer}
\end{figure}

\subsection{Action Layer Decomposition}

A critical insight is that the Action Layer comprises two subcomponents:

\textbf{Contact Layer (Routing/Policy):}
\begin{itemize}
    \item Capability routing
    \item Provider selection
    \item Policy evaluation
    \item Environment selection
\end{itemize}

\textbf{Execution Runtime (Invocation):}
\begin{itemize}
    \item API invocation
    \item Tool execution
    \item Error handling
    \item Result normalization
\end{itemize}

This decomposition maintains clear separation of concerns. The Contact Layer makes \emph{decisions} (which provider, which environment), while the execution runtime performs \emph{actions} (invoke API, return result).

When mapped to the logical model $S = (A, P, E)$, the Contact Layer corresponds to $P$, and the Execution Runtime corresponds to $E$.

\subsection{Defining Statement}

The execution layer is defined as follows:

\begin{quote}
\emph{The execution layer is the action interface between AI cognition and the real-world capability ecosystem.}
\end{quote}

This definition captures the essential role: mediating between what agents decide to do and how those decisions are realized through capability execution.
